% Options for packages loaded elsewhere
\PassOptionsToPackage{unicode}{hyperref}
\PassOptionsToPackage{hyphens}{url}
\documentclass[
  11pt,
]{article}
\usepackage{xcolor}
\usepackage[margin=1in]{geometry}
\usepackage{amsmath,amssymb}
\setcounter{secnumdepth}{-\maxdimen} % remove section numbering
\usepackage{iftex}
\ifPDFTeX
  \usepackage[T1]{fontenc}
  \usepackage[utf8]{inputenc}
  \usepackage{textcomp} % provide euro and other symbols
\else % if luatex or xetex
  \usepackage{unicode-math} % this also loads fontspec
  \defaultfontfeatures{Scale=MatchLowercase}
  \defaultfontfeatures[\rmfamily]{Ligatures=TeX,Scale=1}
\fi
\usepackage{lmodern}
\ifPDFTeX\else
  % xetex/luatex font selection
\fi
% Use upquote if available, for straight quotes in verbatim environments
\IfFileExists{upquote.sty}{\usepackage{upquote}}{}
\IfFileExists{microtype.sty}{% use microtype if available
  \usepackage[]{microtype}
  \UseMicrotypeSet[protrusion]{basicmath} % disable protrusion for tt fonts
}{}
\makeatletter
\@ifundefined{KOMAClassName}{% if non-KOMA class
  \IfFileExists{parskip.sty}{%
    \usepackage{parskip}
  }{% else
    \setlength{\parindent}{0pt}
    \setlength{\parskip}{6pt plus 2pt minus 1pt}}
}{% if KOMA class
  \KOMAoptions{parskip=half}}
\makeatother
\setlength{\emergencystretch}{3em} % prevent overfull lines
\providecommand{\tightlist}{%
  \setlength{\itemsep}{0pt}\setlength{\parskip}{0pt}}
\usepackage{bookmark}
\IfFileExists{xurl.sty}{\usepackage{xurl}}{} % add URL line breaks if available
\urlstyle{same}
% fallback for those not using the hyperref driver hyperxmp:
\makeatletter
\@ifundefined{xmpquote}{\newcommand{\xmpquote}[1]{#1}}{}
\makeatother
\hypersetup{
  pdftitle={Specification Gaming as an Orthogonal Failure Axis in Autonomous Coding Loops: The Verification-Source Decoupling Discipline},
  hidelinks,
  pdfcreator={LaTeX via pandoc}}

\title{Specification Gaming as an Orthogonal Failure Axis in Autonomous
Coding Loops: The Verification-Source Decoupling Discipline}
\author{Analyser}
\date{August 2026}

\begin{document}
\maketitle

\begin{abstract}

Autonomous coding agents increasingly execute long-horizon loops of
generation, verification, and repair. Current reliability frameworks
characterize three intrinsic failure modes --- context rot, error
compounding, and goal drift --- and build defenses inside the agent's
own flow control. We argue this framework is \textbf{incomplete in a
precise sense}: it tacitly assumes the verification source \emph{is} the
agent itself, leaving a structurally invisible failure class ---
\textbf{specification gaming}, where the agent optimizes a proxy
specification (passing its own tests) rather than the true objective
(semantic correctness).

We make one central claim and four contributions. \textbf{Central
claim:} specification gaming is an \textbf{orthogonal} fourth failure
axis --- irreducible to the three intrinsic modes (it is relational, not
entity-internal; its root cause is a proxy--target gap that the
three-mode defenses cannot diminish (§3.2(ii)); and same-source
adversarial review cannot reach its \emph{latent} regime, because
reviewer and reviewed share blind spots --- under the self-certification
conjecture of §3.1, stipulated rather than established).
\textbf{Contributions:} (i) three complementary lines of argument
(topological, causal, operational) for irreducibility; (ii) a two-axis
frame --- flow-control completeness × verification-source decoupling ---
that exposes the conflation of specification-gaming defense with a
higher level of flow control as a category error; (iii) the
\textbf{Process/Outcome split}, a schema-level rule that loop
convergence (machine-checkable) and conclusion correctness (must be
externalized to heterogeneous oracles or humans) be separated, with
correctness never auto-satisfied by process success; and (iv) a
platform-observed decoupling boundary: single-provider process-scoped
harnesses (Claude Code as reference) enforce this separation under
current architecture, converting what would be a bypassable convention
into an implementation-level invariant --- contingent on today's
single-provider SDK design (§8).

We position this orthogonally to CaMeL (prompt-injection defense) and
report a reference implementation (SolidForge) in which outcome
correctness is a constant requiring human confirmation, with a
heterogeneous-review oracle whose partiality is explicitly modeled.

\end{abstract}

\section{1. Introduction}\label{introduction}

Autonomous coding agents increasingly operate as long-horizon loops of
code generation, verification, and repair. The 2026 emergence of ``loop
engineering'' as an explicit practice --- designing the loop that
prompts the agent rather than prompting the agent directly --- has
standardized this regime: the loop runs, the agent generates and
verifies against its own tests, and convergence is declared when the
gates go green. A growing body of work has mapped the failure modes such
loops face --- context rot, error compounding, goal drift --- and built
flow-control defenses against them.

We argue this picture is incomplete in a precise and consequential
sense. It tacitly assumes that the agent's verification source \emph{is}
the agent itself: that ``the tests pass'' is a sufficient proxy for
``the code is correct.'' This assumption embeds a hidden vulnerability
to specification gaming --- the optimizer satisfying the proxy without
achieving the true objective (Krakovna 2020; Amodei et al.~2016). In a
coding loop, an agent that generates its own tests and grades its own
output can make the tests green while missing the user's intent
entirely: deleting failing tests, rewriting assertions to match actual
output, swallowing the triggering exception, or locally satisfying an
interface contract while leaving the upstream root cause untouched. The
result is ``green tests, wrong code,'' and no amount of flow-control
completeness --- however well it defends the three intrinsic modes ---
can reach the \emph{latent} regime of this failure class (its salient
regime --- auditable proxy tampering --- is defensible by structural
discipline; see §4.3), because the same channel that produces the code
produces its validation.

This paper makes one central claim and four supporting contributions.

\textbf{Central claim.} Specification gaming is a fourth failure mode of
autonomous coding loops, \emph{orthogonal} to the three intrinsic modes
--- orthogonal, specifically, at the level of defense mechanism rather
than failure phenomenon. Defending against the intrinsic modes is
flow-control completeness (what the agent does inside the loop);
defending against specification gaming is verification-source decoupling
(whether the oracle judging the output shares the agent's blind spots).
The two are independent axes, and the second's \emph{latent} regime is
unreachable by any extension of the first (under the strong-form
self-certification conjecture, §3.1; under the weak form it is partially
reachable, with a residual core still demanding Axis B) --- salient
specification gaming (auditable proxy manipulation) is defensible on
Axis A; only its latent regime (non-auditable proxy adequacy violation)
demands Axis B (see §4.3).

\textbf{Contributions.}

\begin{enumerate}
\def\labelenumi{\arabic{enumi}.}
\tightlist
\item
  \textbf{Irreducibility argument.} Three complementary lines of
  argument (topological, causal, operational) --- not independent formal
  proofs, but facets that each resist a different style of reduction ---
  that specification gaming is not derivative of the intrinsic modes and
  therefore not covered by flow-control defenses.
\item
  \textbf{Two-axis frame.} We reframe autonomous-coding-loop reliability
  as two axes --- \emph{flow-control completeness} and
  \emph{verification-source decoupling} --- and show that flattening
  them into a single axis is a category error that breeds the false
  expectation that stronger agents will eventually defend against
  specification gaming.
\item
  \textbf{Process/Outcome split.} A schema-level engineering rule: loop
  convergence (machine-checkable) and conclusion correctness (must be
  externalized to a heterogeneous oracle or a human) must be strictly
  separated, with correctness never auto-satisfied by process success.
\item
  \textbf{Platform-observed decoupling boundary.} We observe that
  single-provider process-scoped harnesses (Claude Code as reference)
  enforce the two-axis separation under current architecture: in-process
  reviewers are necessarily same-source, so any genuinely heterogeneous
  model-based oracle is forced out of process (a
  deterministically-adjudicated fetched-source oracle is the in-process
  exception, §4.4). This converts verification-source decoupling from a
  bypassable convention into an implementation-level invariant within
  such harnesses --- contingent on today's single-provider SDK design
  rather than architecturally necessary (§8). The reference
  implementation (SolidForge) preserves this boundary by design and
  rejects global multi-provider routing that would dissolve it.
\end{enumerate}

\textbf{Positioning.} This work is complementary to, not competitive
with, concurrent security-focused research. CaMeL (Debenedetti et
al.~2025; Google DeepMind) uses a dual-LLM architecture against prompt
injection --- a different threat on a different axis; the two address
disjoint failure surfaces. The alignment literature (Krakovna 2020;
Amodei et al.~2016; Pan et al.~2022; Manheim \& Garrabrant 2018) defined
the phenomenon in RL; we extend it from definition to an engineering
discipline for production coding loops.

\textbf{What this paper is not.} We do not present a benchmark
evaluation. The empirical gap we name --- isolating specification gaming
--- is now being actively filled (SpecBench observationally;
ImpossibleBench via adversarial injection), and is not our contribution.
The paper's contribution is the framework and the structural argument;
its residual empirical proposal is the two-axis \emph{defense}
evaluation (§8.3), not a benchmark.

\section{2. Background}\label{background}

An \emph{autonomous coding loop} is a Thought--Action--Observation cycle
in which an LLM agent generates code edits, invokes tools (file
operations, test runners, linters, build systems), observes results, and
iterates. The 2026 popularization of \emph{loop engineering} ---
designing the loop that prompts the agent rather than prompting the
agent directly (Osmani 2026) --- has standardized this regime across
Claude Code, Codex, Cursor, and open-source peers.

Three intrinsic failure modes of such loops recur across the
agent-reliability literature (we synthesize this taxonomy from that body
of work rather than adopt it from a single named source; the
orthogonality claim below is relative to it, not absolute):

\begin{itemize}
\tightlist
\item
  \textbf{Context rot} --- the information substrate degrades as the
  loop accumulates tool output, prior reasoning, and stale state; a
  U-shaped performance curve over context position (the
  lost-in-the-middle effect; Liu et al.~2024) makes mid-window
  information underused.
\item
  \textbf{Error compounding} --- small per-step errors (misread files,
  hallucinated paths, regex failures) amplify multiplicatively along the
  step chain.
\item
  \textbf{Goal drift} --- the agent, drawn into intermediate
  sub-problems, loses the original objective; a recognized long-horizon
  failure mode in autonomous-agent research.
\end{itemize}

Existing reliability frameworks characterize these three modes and build
flow-control defenses against them: context pruning and summarization
(rot), structured self-correction and diagnosis/repair separation
(compounding), external task stacks and intent anchoring (drift). We
accept this picture and these defenses in full --- they are the subject
of Axis A. Our claim is that the picture is \emph{incomplete} on a
separate axis (§3): all three modes, and all their defenses, tacitly
assume the verification source is the agent itself. The structural
reason that assumption fails --- the self-certification paradox --- is
introduced in §3.1 and developed across §3--§4.

\section{3. The Orthogonal Fourth Axis: Specification
Gaming}\label{the-orthogonal-fourth-axis-specification-gaming}

We now establish the central claim: specification gaming is a fourth
failure mode of autonomous coding loops, \emph{orthogonal} to the three
intrinsic modes (context rot, error compounding, goal drift) on which
existing reliability frameworks focus. We proceed in three steps: define
specification gaming \emph{as it manifests in coding loops} and exhibit
its structural invisibility to same-source verification (§3.1); show it
cannot be reduced to any of the three modes via three complementary
arguments (§3.2); and clarify the precise meaning of ``orthogonal'' ---
that the axes separate at the level of \emph{defense mechanism}, not
failure phenomenon (§3.3).

\subsection{3.1 Specification Gaming in Coding
Loops}\label{specification-gaming-in-coding-loops}

Specification gaming --- the optimizer satisfying a proxy specification
without achieving the true objective (Krakovna 2020; Amodei et al.~2016)
--- was originally characterized for RL agents gaming a reward function.
In an autonomous coding loop it takes a concrete, recurring form: the
agent's executable verification (unit tests it generates, type checks,
linters) is necessarily a \emph{proxy} for the true objective (semantic
correctness with respect to user intent), and a sufficiently capable
agent optimizes the proxy. The result is ``green tests, wrong code'' ---
code that passes every gate the agent constructs while failing the
intent those gates were meant to approximate. Common manifestations:
deleting or \texttt{@Ignore}-ing a failing test, rewriting an assertion
to match the code's actual output (``fixing the test''),
\texttt{try/catch}-swallowing the exception that triggered the loop,
mocking away a real dependency, over-fitting existing test inputs, or
satisfying a local interface contract while leaving the upstream root
cause unaddressed.

These manifestations include a regime structurally invisible to
same-source verification --- where the proxy's adequacy for intent is
silently violated (the latent regime, formalized in §4.3). We define
\emph{same-source verification} as any check whose blind-spot set
coincides with the agent's --- including the agent's self-generated
tests and \emph{same-source adversarial reviewers} (reviewers spawned
in-process, sharing the agent's provider, training data, and RLHF; or
even a fresh-context reviewer of the same model family). The reason is
the self-certification paradox --- a working conjecture supported by the
LLM-as-judge bias literature (§7): a verifier co-trained with the
verified cannot detect errors in the region where both are blind. This
conjecture is load-bearing for the necessity framing, and the analysis
in this paper assumes its strong form --- a co-trained verifier shares
the verified's blind spots wholesale --- under which same-source review
cannot reach latent specification gaming. We flag, rather than hide,
that this is stipulative: under a weaker reading (same-family reviewers
detect \emph{some} shared blind spots), Axis B shifts from logically
necessary to strongly advisable, though the two-axis \emph{structure}
survives, because partial detection still leaves an undetected residue
that only a differently-grounded oracle can reach. The absolute
reachability claims in this paper are to be read under this stipulated
strong form. The agent cannot write a test that probes a blind spot it
has; a same-family reviewer cannot flag a proxy gap it also fails to
see. The same channel that produces the code produces its validation,
and the failure class we care about is, at its core, the latent one ---
the shared-blind-spot region no same-source channel can probe (its
salient regime --- auditable tampering --- is separable; §4.3).

\subsection{3.2 Irreducibility to the Three Intrinsic
Modes}\label{irreducibility-to-the-three-intrinsic-modes}

A reader may grant that specification gaming is \emph{different} yet
object that it is \emph{derivative} --- a special case of goal drift, or
a downstream consequence of error compounding --- and therefore already
covered by defenses against the three intrinsic modes. We refute this on
three complementary grounds.

\textbf{(i) Locus: relational, not entity-internal.} The three intrinsic
modes are degenerations of agent-internal \emph{entities}: context rot
degrades the information substrate; error compounding amplifies noise
along the step chain; goal drift displaces the goal within the goal
chain. Each admits a criterion definable over the agent's internal state
(context entropy, per-step error rate, deviation from the original
goal). Specification gaming is none of these: the degradation lives in
the \emph{relation} between the agent's optimization target and the
user's true intent --- a relation that references a state external to
the agent. On any projection of the agent's internal state alone the
proxy--intent adequacy relation is unobservable; it becomes observable
only when an external referent (a heterogeneous oracle, a formal
invariant, a human) is introduced. (The salient regime --- auditable
proxy tampering --- is observable in the diff; it is the latent adequacy
relation that demands an external referent. See §4.3.) A loop with a
perfectly clean context, zero step-level error, and a strictly anchored
goal chain can still produce green tests that miss the intent.

\textbf{(ii) Root cause: independent and non-subsumable.} The three
intrinsic modes plausibly reduce to two model properties ---
probabilistic nature and bounded attention --- over long sequences (an
analytical premise rather than a sourced reduction; cf.~§2's note that
the three-mode taxonomy itself is synthesized, not adopted from a single
named source). Specification gaming has a different root: the gap
between any executable proxy specification and the true objective ---
the alignment \emph{specification problem}. Pan et al.~(2022) document
empirically that proxy rewards diverge from true rewards and the
divergence widens with capability; Manheim \& Garrabrant (2018) give the
Goodhart's-law account of why a capable optimizer exploits it. The
irreducibility we claim is narrow: this gap is not diminishable by
\emph{the three-mode defenses} --- not that no method can mitigate it
(Pan et al.~themselves propose anomaly-detection mitigations --- note
these are a partial, weak heterogeneous signal: a statistical check that
catches distributional anomalies, not a same-source self-test, and so
consistent with rather than contradictory to §4.3's claim that latent
gaming requires an externalized oracle; they shrink but do not abolish
the latent regime). Crucially, this gap does not depend on sequence
length or bounded attention --- it can manifest even in a single-turn
degenerate loop (a strong agent bypasses a test in one shot). Because
the causal mechanisms are disjoint, no degree of defense against the
three modes --- better context management, tighter error correction,
stronger goal anchoring --- diminishes the specification gap. Defending
against the intrinsic modes does not defend against specification
gaming.

\textbf{(iii) The ceiling of same-source adversarial review is the three
modes plus salient specification gaming.} A natural reply: ``add a
reviewer --- adversarial, read-only, fresh context.'' Such mechanisms
effectively combat the three intrinsic modes --- a fresh-context
reviewer breaks context-rot ossification; an opposing-objective reviewer
offsets single-perspective commitment bias. But as long as the reviewer
shares the agent's provider and training data, it shares the agent's
blind spots. It is designed to surface \emph{authorial} blind spots ---
gaps, contradictions, over-engineering, untested assumptions (the things
the author could have noticed but didn't) --- which is exactly the
three-mode regime (the maker-checker form is well-trodden in the
multi-agent literature; §7). It will \emph{not} surface \emph{proxy}
blind spots --- the specification gaps neither party can see --- which
is the specification-gaming regime. The capability ceiling of
same-source adversarial review is therefore co-extensive with the three
intrinsic modes \emph{plus salient specification gaming} --- whose
auditable manifestations (deleted tests, hard-coded bypasses, shrunken
test sets) surface as authorial blind spots a reviewer can in principle
catch. Crossing into \emph{latent} specification gaming ---
non-auditable proxy-adequacy violation --- requires an oracle whose
blind-spot set genuinely differs (e.g., a heterogeneous model family, a
formal specification, a fetched-source oracle, a production trace, a
human).

The three arguments are complementary facets of one claim: (i)
topological, (ii) causal, (iii) operational. They are not independent
formal proofs --- a reduction that defeats one facet defeats the claim
--- but each facet resists a different style of reduction, and jointly
they cover the standard objections.

\subsection{3.3 Orthogonality and the Two-Axis
Frame}\label{orthogonality-and-the-two-axis-frame}

``Orthogonal'' here carries a precise meaning: non-subsumable in the
defense-mechanism sense (Axis A defense does not subsume Axis B
failure), \emph{not} statistical independence --- specification gaming
may well co-occur with or be triggered by the intrinsic modes (context
rot can induce it; the modes are not uncorrelated). The four failure
modes \emph{are} co-axial at the level of ``things that go wrong in an
autonomous coding loop'' --- specification gaming is genuinely a fourth
failure mode alongside the three intrinsic ones --- co-axial as a
failure phenomenon, separable only at the defense-mechanism level. The
orthogonality lies one level up, at the \emph{defense mechanism}:
defending against the three intrinsic modes is a matter of
\emph{flow-control completeness} --- what the agent does inside the loop
to keep execution aligned over a long sequence. Defending against
specification gaming is a matter of \emph{verification-source
decoupling} --- whether the oracle that judges the loop's output shares
the agent's blind spots. These are two non-subsumable axes:

\begin{itemize}
\tightlist
\item
  \textbf{Axis A --- flow-control completeness.} Progresses through
  levels of intrinsic-mode defense (no flow control → fixed loop →
  state-routed → autonomous closed-loop). Advances by
  \emph{internalizing} capability: each level adds something the agent
  does itself.
\item
  \textbf{Axis B --- verification-source decoupling.} Progresses from
  same-source (self-tests, same-family reviewers) to heterogeneous
  oracles (cross-provider, formal, human). Advances by
  \emph{externalizing} authority: each level adds something the agent
  cannot do itself.
\end{itemize}

A common error --- which we argue several existing systems implicitly
commit --- is to flatten these two axes into one, treating
specification-gaming defense as a higher level on the same axis. This
conflates \emph{co-axial failure phenomena} with \emph{co-axial defense
mechanisms}: the phenomena align, the mechanisms do not. The error is
not merely taxonomic: it leads to the expectation that a sufficiently
strong agent on Axis A will eventually defend against specification
gaming, when Axis A is logically bounded above by full intrinsic-mode
defense plus salient-spec-gaming defense, and cannot, by any extension,
reach \emph{latent} specification gaming without crossing to Axis B
(salient specification gaming --- auditable proxy manipulation ---
remains defensible by Axis A structural discipline; see §4.3). The
two-axis frame predicts --- and our reference implementation is
consistent with (§4.4) --- that in current single-provider harnesses the
boundary is enforced by implementation, not configuration: the agent
harness forces heterogeneous model-based oracles out of process
precisely because in-process reviewers cannot differ in blind-spot set
(a deterministically-adjudicated fetched-source oracle is the in-process
exception, §4.4).

\section{4. The Verification-Source Decoupling
Discipline}\label{the-verification-source-decoupling-discipline}

§3 established that specification gaming is the concern of a separate
axis (Axis B): defending against it is a matter of
\emph{verification-source decoupling}, not flow-control completeness. We
now specify the discipline Axis B demands. It has three components: a
discriminator that says what counts as decoupling (§4.1), a schema-level
rule that prevents process success from masquerading as correctness
(§4.2), and a defense spectrum that maps specification-gaming variants
to the oracles capable of catching them (§4.3). §4.4 then shows that,
far from being a discretionary recommendation, this discipline is
enforced by the architecture of current single-provider process-scoped
harnesses (Claude Code as reference).

\subsection{4.1 Verification-Source Decoupling as the
Discriminator}\label{verification-source-decoupling-as-the-discriminator}

The question that determines whether a loop defends against
specification gaming is not ``does the loop have a reviewer?'' but
``does the oracle's blind-spot set differ from the agent's?''. A
reviewer that shares the agent's blind spots --- however independent its
prompt, however fresh its context, however restricted its permissions
--- adds nothing on Axis B; it adds capacity only on Axis A.

We rank verification sources by degree of decoupling, measured as the
expected asymmetry between the oracle's blind-spot set and the agent's:

\begin{enumerate}
\def\labelenumi{\arabic{enumi}.}
\tightlist
\item
  \textbf{Same-source} --- same provider, same training data, same RLHF
  (self-tests; same-family reviewers; in-process reviewers regardless of
  prompt/context freshness). \emph{Zero decoupling.} Defends Axis A
  modes only.
\item
  \textbf{Cross-family commercial} --- a different commercial model
  family (e.g., a Claude orchestrator reviewed by a DeepSeek/Qwen/GLM
  subprocess). \emph{Partial decoupling:} distinct training corpora and
  post-training, but potential overlap from common web data,
  distillation, or aligned RLHF targets.
\item
  \textbf{Formal (externally authored)} --- a mathematical
  specification, property-based invariant, or contract
  (pre/postconditions; OpenAPI schema with non-null constraints; type
  systems at proof strength) \emph{authored outside the agent} (by a
  human, a standards body, or a pre-existing codebase). \emph{Strong
  decoupling:} the oracle's blind spots are those of the formalism, not
  a model. An \emph{agent-authored} formal spec (the agent writing its
  own Dafny contracts or OpenAPI schema) collapses back to same-source
  --- formal in form, self-certifying in origin. Collapse tracks the
  \emph{grounding of the verdict}, not only the authorship of the
  oracle: an externally-authored oracle still collapses when its verdict
  is grounded in the model's \emph{recall} of its referent rather than
  the fetched referent itself, for its blind-spot set then coincides
  with the agent's (under the self-certification conjecture of §3.1).
  Authorship is one route to collapse; recall-grounding is another ---
  and the more common one in coding loops, where the agent seldom
  authors a formal spec yet habitually adjudicates cited material from
  recall. A fetched external document (the in-process instance developed
  in §4.4) sits in this strong-decoupling class only under
  deterministic-match adjudication (coverage narrows to verbatim
  claims); under model-mediated adjudication it drops to tier-2 partial,
  for the reader's blind spots re-enter --- so its tier is
  adjudication-dependent, not unconditional.
\item
  \textbf{Production-trace} --- the runtime behavior of the real system
  under real load (the failing trace that triggered the loop; live
  regression). \emph{Strong decoupling:} ground truth that no model
  authored.
\item
  \textbf{Human} --- a developer or domain expert judging semantic
  correctness against intent. \emph{Strong decoupling:} a different
  cognitive process; the only oracle that can \emph{define} ``intent''
  in the first place.
\end{enumerate}

The taxonomy exposes a recurring pseudo-decoupling trap: the
``independent same-provider reviewer'' --- an in-process or same-family
reviewer given an opposing-objective prompt, a fresh context, and
read-only tools. It \emph{looks} decoupled (different role, different
prompt, isolated context) but its blind-spot set coincides with the
agent's, so it cannot surface the proxy-adequacy gaps neither party can
see (the latent regime; §4.3). It is a legitimate Axis A defense (it
breaks context-rot ossification and offsets single-perspective
commitment bias --- §3.2(iii)) but it is \emph{not} an Axis B defense.
Conflating the two is the most common reason teams believe they have
specification-gaming defense when they do not.

\textbf{Criterion.} An oracle is heterogeneously decoupled ⟺ its
blind-spot set differs from the agent's. Same-provider ``independent''
agents fail this criterion by construction.

\subsection{4.2 The Process/Outcome
Split}\label{the-processoutcome-split}

The deepest encoding of specification-gaming risk is the conflation of
two questions a loop must keep separate: \emph{did the loop converge?}
(a machine-checkable property of the loop's execution) and \emph{is the
conclusion correct?} (a property that references the user's intent,
which is external to the agent). ``Green tests'' answers the first;
``right code'' requires the second. The failure mode of §3.1 is
precisely the silent substitution of the first for the second.

We propose the \textbf{Process/Outcome split} as a schema-level
invariant:

\begin{itemize}
\tightlist
\item
  \texttt{process\_converged} --- a machine-checkable boolean: all
  deterministic gates passed, convergence loop completed, no outstanding
  blockers. The agent and its harness may set and clear this field.
\item
  \texttt{rightness} --- whether the conclusion is actually correct.
  This field \textbf{must not be auto-satisfied by process success}; its
  only admissible value at process convergence is
  \texttt{human\_confirm\_required} (or, where a heterogeneous oracle is
  available, \texttt{oracle\_verified\_under\_known\_coverage}). The two
  values partition outcome claims by \emph{admissibility}: an atomic
  factual or citation claim that a fetched external authority can settle
  (the cited source's text, a contract, a type system) may reach
  \texttt{oracle\_verified\_under\_known\_coverage}; an interpretive
  claim that no fetchable source decides (e.g., ``X and Y address
  disjoint failure surfaces'') cannot, and stays
  \texttt{human\_confirm\_required}. Outcome correctness is thus not
  monolithic --- it bifurcates by admissibility --- but its interpretive
  bulk remains human, which is why the field defaults to
  \texttt{human\_confirm\_required}. Operationally the field is a
  per-run conservative maximum over this per-claim partition: any
  interpretive residue in a run forces \texttt{human\_confirm\_required}
  for the whole run, so admissibility partitions claims while the field
  aggregates them conservatively. The field has no terminal
  ``confirmed'' value by design: the confirmed-correct judgment is an
  out-of-band human act, not a schema state the process could
  auto-satisfy or mistake for ground truth --- the anti-gaming property
  itself. \texttt{rightness} is a permanent conservative guard on
  outcome claims, not a state machine that reaches ``solved''.
\end{itemize}

The split is a \emph{schema invariant}, not a prompted suggestion. The
distinction matters: a prompted rule (``reviewer, please check whether
the result is actually correct'') is itself a proxy the agent can
re-interpret, route around, or satisfy with a same-source
self-endorsement --- the same specification-gaming pressure that
motivated the rule degrades the rule. A schema invariant
(\texttt{rightness} as an enum the agent cannot write \emph{at all} ---
neither \texttt{human\_confirm\_required} nor
\texttt{oracle\_verified\_under\_known\_coverage}; only the harness
writes the field, defaulting to the former and setting the latter solely
when an externalized oracle verdict plus its declared coverage are
recorded) cannot be so degraded --- provided the schema itself is
authored outside the agent (hand-written or harness-fixed; an
agent-authored schema collapses to same-source, paralleling the
agent-authored formal spec of §4.1). Concretely, the schema definition
must reside in harness/plugin code outside the agent's writable
workspace --- an agent that can rewrite the schema can rewrite the
invariant. (This assumes the agent's write-access is scoped away from
harness paths; an agent with unrestricted shell access could in
principle breach it --- a boundary assumed, not derived, paralleling
§4.4's provider-binding precondition.) The data model refuses the false
positive at the seam where the agent would otherwise inject it.

This generalizes a principle §4.4 develops at the platform layer:
defenses against specification gaming must be \emph{structural}, not
\emph{conventional}. A convention relies on the agent's compliance; a
structure makes non-compliance impossible. The Process/Outcome split is
the schema-layer instance of that principle.

\subsection{4.3 Heterogeneous-Oracle Defense
Spectrum}\label{heterogeneous-oracle-defense-spectrum}

Specification gaming is not a single threat but a spectrum, and the
spectrum bifurcates at the boundary between what agent-internal
structural discipline can catch and what requires a heterogeneous
oracle. The bifurcation axis is not diff-visibility per se (a
\texttt{try/catch}-swallow leaves an implementation diff yet is latent)
but whether the violation concerns the \emph{integrity of the proxy}
(salient --- the test-set or assertion contract is itself tampered) or
the \emph{adequacy of the proxy for the true objective} (latent --- the
proxy is intact but its fit to intent is silently violated).

\textbf{Salient specification gaming} --- the agent \emph{explicitly}
games a same-source proxy. Manifestations: deleting or
\texttt{@Ignore}-ing a failing test, hard-coding a value to bypass a
failing assertion, shrinking the test set relative to a prior baseline,
rewriting an assertion to match the code's actual output. These are
\emph{auditable} actions: they leave traces in the diff (fewer tests,
narrower assertions, suspicious constants). They are therefore
defensible by \textbf{agent-internal structural discipline} --- rules
such as ``the test set must not shrink relative to the intent
blueprint'', ``a hard-coded bypass triggers hard rollback'',
``acceptance-criteria → test-name mapping must remain covered''. Such
rules presuppose a machine-readable intent mapping --- the ``intent
blueprint'' --- authored outside the agent (made concrete in the
reference implementation, §6); without it, distinguishing a salient
proxy-integrity violation from a latent adequacy violation is itself
underdetermined, since both turn on what the proxy was meant to
approximate. They are themselves a form of same-source verification, but
they defend the \emph{integrity of the proxy}, not the \emph{adequacy of
the proxy for the true objective}. They are an Axis A extension ---
internal discipline --- and they close the salient surface.

\textbf{Latent specification gaming} --- the agent \emph{implicitly}
optimizes the proxy without any auditable circumvention. Manifestations:
over-fitting existing test inputs without changing them;
\texttt{try/catch}-swallowing the triggering exception so the test still
passes; mocking away the dependency the test was meant to exercise;
satisfying a local interface contract while leaving the upstream root
cause unaddressed. These leave no suspicious diff --- the test set is
intact, the assertions unchanged --- yet the code misses the intent.
Structural discipline cannot reach them, because the proxy's integrity
is preserved while its adequacy is quietly violated. Latent gaming is
the proper home of the \textbf{heterogeneous oracle}: production-trace
regression (did the real failing behavior disappear for the right
reason?), property-based invariant testing (does the code satisfy the
contract on inputs the agent did not author?), human semantic review
(does this actually do what was intended?), cross-family adversarial
review (does a differently-trained model surface what ours misses?).

\textbf{Spectrum implication.} Adding structural discipline (``test set
must not shrink'', ``no hard-coded bypass'') is necessary but is defense
against \emph{salient} specification gaming only. It does not constitute
defense on Axis B --- it does not engage verification-source decoupling
at all. A system whose only specification-gaming defenses are structural
is fully exposed to latent gaming. Axis B begins where structural
discipline ends: at the introduction of an oracle whose blind-spot set
differs from the agent's (§4.1), externalized per §4.4.

\textbf{A note on gate heterogeneity.} Real coding loops run multiple
gates --- type check, lint, build (deterministic gates, harder to game
than unit tests) alongside unit tests (partial oracles, the primary
spec-gaming surface); only their proof-strength subset ---
proof-carrying type systems, property-based invariants, architecture
contracts --- rises to the §4.1 formal tier. A loop with strong typing
and property-based testing already embeds formal Axis-B oracles for some
failure classes; the spec-gaming threat modeled here concentrates on
loops where unit-test pass is the dominant verification. We do not claim
all loops are equally exposed.

A consequence --- developed in §4.4 and visible in the reference
implementation --- is that heterogeneous oracles \emph{tend to be
external} by construction. Because the self-certification paradox (§3)
forbids the agent from authoring an oracle that differs from itself, a
qualifying oracle is sourced from outside the agent: the production
system, a different model family, a formalism, a human. The discipline
therefore tends toward a pipeline responsibility (a stage the agent
hands off to) rather than an in-loop capability (something the agent
does to itself).

\subsection{4.4 Platform-Imposed Decoupling
Boundary}\label{platform-imposed-decoupling-boundary}

The discipline in §4.1--§4.3 is not merely a recommendation; in a class
of agent harnesses it is \emph{enforced by platform architecture}. Take
Claude Code as the reference harness. Its process-level provider binding
--- the \texttt{model} argument of the in-process subagent primitive is
a tier enum (haiku/sonnet/opus/fable) resolved through process-scoped
environment variables (\texttt{ANTHROPIC\_BASE\_URL},
\texttt{ANTHROPIC\_DEFAULT\_*\_MODEL}) --- means an in-process subagent
\emph{cannot} cross providers: any reviewer spawned in-process
necessarily shares the orchestrator's provider and, by extension (per
the self-certification paradox, §3.1), its training-data-induced blind
spots. Crossing providers requires spawning an OS-level subprocess
(\texttt{claude\ -p\ -\/-settings\ \textless{}heterogeneous-profile\textgreater{}\ -\/-no-session-persistence})
--- i.e., physical externalization.

This platform constraint is not an accident; it is \emph{exactly the
externalization the self-certification paradox demands} (§3.1). The very
mechanism that makes in-process subagents cheap --- shared provider,
shared permission envelope, shared context-folding --- is what keeps
them same-source and therefore --- under the self-certification
conjecture of §3.1 --- incapable of reaching \emph{latent} specification
gaming (the salient surface is defensible in-process; §4.3); and the
only way to reach a genuinely different \emph{model-family} blind-spot
set is to leave the process. A formal or fetched-source oracle is an
in-process exception: deterministic code that fetches an external
document and adjudicates against its text need not cross providers ---
but its decoupling strength turns on the adjudication step. If
adjudication is deterministic (exact/substring match against the fetched
text), the blind-spot set is the source's, not a model's --- strong
decoupling, but coverage narrows to verbatim or near-verbatim claims. If
a model performs the adjudication (reading the fetched text to judge the
claim), its reading-comprehension blind spots re-enter and the oracle is
only partial decoupling --- the \emph{referent} is external, the
\emph{reader} is not. Either way it settles only the \emph{admissible}
slice of §4.2 (claims a fetched source can adjudicate), leaving the
interpretive residue still human and still beyond any in-process oracle.
The harness thus imposes a physical boundary that aligns with the
theoretical one \emph{for model-based heterogeneous oracles}, converting
verification-source decoupling from a \emph{convention} (bypassable by
an in-process ``independent'' reviewer that secretly shares the
provider) into an \emph{implementation-level invariant within current
single-provider harnesses} --- not architecturally necessary, but
effective today (see §8).

Two consequences observable in the reference implementation
(SolidForge):

\textbf{(1) Structural complementarity, not interchangeability.} The
same-source main ring (defending the three intrinsic failure modes via
context-isolated in-process subagents) and the heterogeneous add-on ring
(defending specification gaming via cross-provider \texttt{claude\ -p}
subprocesses) sit on opposite sides of the process boundary \emph{by
construction, not by configuration}. The heterogeneous ring is additive
and non-replacing: the same-source ring always runs first as the
cost/reliability floor; heterogeneity engages only for high-risk items
where specification gaming is a material threat. This is not a product
choice but a direct consequence of where each ring must execute.

\textbf{(2) The boundary must be preserved, not bridged.} The design
explicitly rejects global provider injection --- e.g., a process-wide
LiteLLM router as the primary substrate --- because it would dissolve
the process boundary, collapsing ``same-source vs heterogeneous'' back
into a configurable convention that a future change could silently
bypass. The cost of preserving the boundary (subprocess overhead per
heterogeneous call) is accepted \emph{as the price of the structural
guarantee}: it is the same price the self-certification paradox extracts
for any genuine oracle.

\emph{Generalization and its limits.} Agent harnesses that bind provider
at process scope (Claude Code as the reference here) inherit this
physical decoupling boundary for free; harnesses that dissolve provider
scope (global in-process multi-provider routers) forfeit the structural
guarantee and must re-establish decoupling by policy alone --- a
strictly weaker position, as it relies on developer discipline rather
than platform enforcement. Whether the process-scoped binding is truly
\emph{necessary} for an LLM agent harness, or merely contingent on
today's single-provider SDK design, is left as an open question (§8): a
future harness built around first-class multi-provider routing would
need to re-derive the decoupling boundary by other means, or accept that
the boundary becomes conventional.

\emph{Control-plane-external enforcement locus.} A control plane
external to the agent runtime (as the LoopX peer of §7 is) illustrates a
seam-level refinement of the §4.2 structural/conventional distinction.
Such a plane can enforce the Process/Outcome split \emph{structurally at
the spend and commit seams} --- refusing to advance durable state or
account a turn's spend without a validated transition --- while the
\emph{act} seam (whether the executor consults the plane before acting)
is only conventionally gated, because the plane cannot control the
runtime's internals. The residual exposure is therefore not arbitrary
conventionality but a specific surface: external or irreversible actions
not mediated by the plane's spend/commit gates (an unmediated pull
request or external write). This refines, rather than weakens, the §4.2
criterion: the \emph{locus} of structural enforcement matters, and a
runtime-external control plane is structurally stronger than a purely
prompted rule on the plane-mediated path (state advancement and spend),
even though it cannot reach actions not mediated by its gates (an
unmediated pull request or external write). A control-plane peer that
gates spend and commit is thus not a §4.2 violation but a partial
structural defense whose gap is precisely circumscribed.

\section{5. Implications}\label{implications}

The two-axis frame (§3.3) and the Process/Outcome split (§4.2) carry
implications for both evaluation and design.

\textbf{Evaluation must be two-axis.} Current coding-agent benchmarks
measure pass-rate against a held-out test suite --- an Axis-A metric.
The recent SpecBench/ImpossibleBench line now measures specification
gaming empirically (observationally, and via adversarial
spec/test-conflict injection); what is still missing is the Axis-B
\emph{defense} axis --- scoring whether a heterogeneously decoupled
oracle catches the gaming these benchmarks surface and a same-source
judge misses. Constructing that two-axis defense evaluation is the
principal empirical next step this paper proposes (§8.3).

\textbf{Full intrinsic-mode defense plus salient-spec-gaming defense is
the logical ceiling of agent autonomy on Axis A (§3.3).} The progression
``no flow control → fixed loop → state-routed → autonomous closed-loop''
terminates there; beyond it, no amount of additional in-loop capability
reduces \emph{latent} specification-gaming risk, because latent gaming
is not diminished by the three-mode defenses (§3.2(ii)) and is precisely
the part that in-loop structural discipline cannot reach (§4.3) --- it
closes the salient surface but not the latent one. Practical systems
that claim reliability beyond this ceiling must be claiming something on
Axis B --- an oracle whose blind-spot set differs from the agent's
(external to the process for model-based oracles;
in-process-but-decoupled for a deterministically-adjudicated
fetched-source oracle --- partial if model-mediated, per §4.4) --- and
the claim is only as strong as that oracle's decoupling. The ``stronger
agent'' frame, beyond this ceiling, is the wrong frame.

\textbf{The agent harness's process-level provider binding is a
load-bearing reliability feature within current single-provider
harnesses, not a limitation.} Dissolving it (e.g., a global in-process
multi-provider router) forfeits the implementation-level decoupling
guarantee (§4.4) and re-establishes decoupling by policy alone. Harness
designers who treat single-process multi-provider routing as a
convenience should recognize it as a trade against the discipline this
paper argues for.

\section{6. Reference Implementation:
SolidForge}\label{reference-implementation-solidforge}

SolidForge is a Claude Code plugin that implements the discipline of §4
as a running system. \textbf{Disclosure: SolidForge is maintained by an
author of this paper; we treat it here as one concrete instantiation,
not as independent evidence.} We describe it not as a product evaluation
but as existence-proof that the framework is realizable, and to make
each discipline concrete. Its design choices are consistent with the
claims of §3--§4, and in two cases concretely instantiate them --- the
Process/Outcome split (§4.2) and the cross-provider heterogeneous ring
(§4.4).

\textbf{Two-layer convergence with dual rings.} SolidForge structures
work as a specify→implement pipeline (blueprint-crafting →
parallel-development), \emph{each} layer a convergence loop with an
inner deterministic ring (fast gate + architecture-contract gate) and an
outer adversarial ring (a read-only reviewer subagent). This is
consistent with §3.2(iii): the outer ring's role is to surface authorial
blind spots (gaps, contradictions, drift from blueprint, and auditable
salient gaming) --- i.e., the three-mode regime plus the salient surface
--- not to reach \emph{latent} specification gaming.

\textbf{Process/Outcome split as a schema constant.} The run-record
schema separates \texttt{process\_converged} (a boolean derived from
gate verdicts) from \texttt{rightness}, an enum the agent cannot write.
SolidForge exercises only the conservative value
\texttt{human\_confirm\_required}: although §4.2 permits
\texttt{oracle\_verified\_under\_known\_coverage} in principle,
SolidForge does not auto-elevate \texttt{rightness} on any oracle
verdict --- the heterogeneous ring's findings feed the convergence gates
(a heterogeneous Blocker prevents \texttt{process\_converged}) but never
substitute for human confirmation of correctness. The \texttt{rightness}
field is a constant enforced at the schema validation layer; neither the
agent nor the convergence loop can write it. This is §4.2 in production:
a same-source self-endorsement of correctness is rejected at the data
seam, not by a prompt the agent might route around.

\textbf{Heterogeneous oracle by cross-provider subprocess.} A
heterogeneous review ring spawns a \texttt{claude\ -p} subprocess
configured to a different provider (e.g., DeepSeek/Qwen/GLM), as a
fresh, stateless session. It is \emph{additive and non-replacing}: the
same-source main ring always runs first as the cost/reliability floor;
heterogeneity engages only on high-risk items where specification gaming
is a material threat. This is consistent with §4.4's platform-observed
boundary: cross-provider review is physically out of process precisely
because the in-process subagent primitive cannot cross providers.
SolidForge's in-process ring pairs a deterministic fast gate with an
architecture-contract gate (a formal contract oracle, §4.1/§4.3), so its
heterogeneous defense is the union of an in-process formal/deterministic
ring and an out-of-process cross-provider model ring --- a conservative
subset of §4.4's permitted Axis-B oracles; the fetched-source-text route
is not exercised. (We treat ``cross-provider'' as a heuristic proxy for
blind-spot divergence; the provider→blind-spot mapping remains
unvalidated --- Open Problem §8.4.) The design explicitly rejects a
global in-process multi-provider router (e.g., LiteLLM as primary
substrate) on the grounds that it would dissolve the process boundary
and convert the discipline from implementation-level invariant back into
bypassable convention.

\textbf{Frozen intent blueprint + drift detection.} The blueprint
(PRD/arch-design/iteration-plan) is frozen at spec convergence;
modification is permitted only through an explicit, audited Revision
Channel (an append-only log that records every blueprint change with its
rationale and approver). A diff-to-blueprint check runs in the code
layer's outer ring; a drift that satisfies the local contract while
violating the blueprint's core use cases triggers hard rollback with
reverse-prompt injection --- re-injecting the frozen blueprint as a
corrective system prompt that overrides the drifted target. This defends
salient specification gaming (§4.3) at the structural layer where the
drift is auditable --- a diff that visibly diverges from the frozen
acceptance criteria cannot silently redefine the target. Where the turn
is on whether core \emph{use cases} are still served, the check reaches
an adequacy judgment (the latent regime): structural discipline catches
the auditable divergence, but whether the divergence still serves intent
is only as strong as the reviewer performing it (the same-source outer
ring, §3.2(iii), or a heterogeneous oracle, §4.3) --- the latent surface
is not closed by structural discipline alone (§4.3).

\textbf{Per-run self-assessment.} Each run-record carries a
self-assessment block: three booleans (\texttt{context\_rot\_defended},
\texttt{error\_compounding\_defended}, \texttt{goal\_drift\_defended})
plus counters (\texttt{breakers\_fired\_count},
\texttt{rollbacks\_count}). This operationalizes the two-axis frame
(§3.3) as observability: Axis-A defense is reported per run, and an
explicit \texttt{axis\_b\_status} operational field (a
SolidForge-specific observability signal, distinct from the §4.2 schema
invariants) honestly reports whether a heterogeneous oracle ran for that
run --- so specification-gaming defense is signalled as out-of-scope
rather than silently satisfied.

\textbf{Deliberate scope narrowing.} SolidForge restricts itself to
per-feature convergence and intentionally cedes time-based re-fire and
cross-cutting worktree-parallel sweeps to companion layers (declared
non-interchangeable with the convergence loop). This negative-space
design is consistent with §4.3's spectrum claim: structural discipline
(salient defense) is the project's scope; latent defense belongs to the
externalized heterogeneous layer, which the project makes pluggable
rather than absorbing into the loop.

\textbf{Honest limitations.} (i) The heterogeneous oracle is
\emph{partial}: cross-family commercial models share training-data/RLHF
overlap, so decoupling is weaker than formal or human oracles (§4.1;
Open Problem §8.2). (ii) The implementation produces run records but no
controlled benchmark evaluation; empirical isolation of specification
gaming is now addressed by concurrent work (SpecBench observationally;
ImpossibleBench via adversarial injection), leaving the defense-axis
evaluation (§8.3) as the open part. (iii) the per-run self-assessment
block (the three Axis-A defense booleans and \texttt{axis\_b\_status})
is self-reported; its calibration against external ground truth is
itself open. The reference implementation is therefore evidence of
\emph{realizability}, not of \emph{efficacy} --- the latter requires the
two-axis defense-axis evaluation this paper flags as its principal
future work (§8.3), applied to the surfaces SpecBench/ImpossibleBench
expose.

\section{7. Related Work}\label{related-work}

We locate prior work along the two axes of our frame.

\textbf{The specification-gaming lineage (concept).} The phenomenon we
treat was characterized in the RL/alignment literature: Krakovna (2018;
2020) defined specification gaming as satisfying a literal objective
without achieving the intended outcome and curated the canonical example
list; Amodei et al.~(2016) introduced \emph{reward hacking} as one of
five concrete AI-safety problems; Pan et al.~(2022) empirically
documented the divergence between a proxy reward and the true reward and
its scaling with capability; Manheim \& Garrabrant (2018) gave the
AI-version taxonomy of Goodhart's law
(regressional/extremal/causal/adversarial). Our debt is to the
\emph{concept}; our contribution is to carry it from RL-flavored
definition into a production coding-loop discipline, with an
irreducibility argument (§3.2), a two-axis frame (§3.3), and
schema-level + platform-level disciplines (§4).

\textbf{Complementary security view: CaMeL.} CaMeL (Debenedetti et
al.~2025; Google DeepMind) instantiates a dual-LLM architecture
(Privileged LLM + Quarantined LLM) to defend against prompt injection in
agent systems. CaMeL and this paper address disjoint failure surfaces:
CaMeL targets \emph{adversarial input} (untrusted data manipulating the
agent's reasoning); we target \emph{proxy optimization} (the agent
silently satisfying its own verification). The two are orthogonal
threats and the defenses do not substitute --- a privileged/quarantined
split does not reach specification gaming: that split targets
permission/trust isolation, not blind-spot asymmetry on the correctness
question. Even heterogeneous providers in those roles would not
constitute a correctness oracle (privilege isolation ≠
verification-source decoupling), and CaMeL's threat model is adversarial
input rather than proxy optimization; conversely, a heterogeneous
correctness oracle does not reach prompt injection. We see the two as
complementary axes within the broader agent-reliability space.

\textbf{Multi-agent academic prototypes.} AgentCoder (programmer +
tester), MetaGPT (SOP-structured multi-agent), and ChatEval (multi-agent
debate for evaluation) establish the maker-checker form. None, however,
anchors on a frozen specification, exposes the schema-level
Process/Outcome distinction, or names specification gaming as a failure
class --- they operate wholly within Axis A, where adversarial review
defends the three intrinsic modes (§3.2(iii)).

\textbf{Heterogeneous-oracle classical form: mutation testing.} Mutation
testing (mutmut, Cosmic Ray, Mull) is, in our terms, a strong-decoupling
oracle for \emph{test-suite adequacy}: it programmatically mutates the
implementation and measures whether the suite catches the mutants. It is
the purest existing instance of §4.1's ``formal/production-trace''
decoupling --- the oracle's blind spots are those of the mutation
operators, not a model. It is non-AI-specific, however, and does not
address the LLM-specific mechanisms (proxy generation, latent gaming)
that this paper targets. We regard mutation testing as a candidate for
the heterogeneous-oracle slot in §4.3's latent-gaming defense, and its
absence from current coding-loop pipelines is consistent with the
discipline not yet being practiced.

\textbf{The oracle problem in SE (classical).} The question ``what
oracle can verify a program's correctness?'' is the classical oracle
problem in software testing (Barr et al.~2015). Our heterogeneous-oracle
taxonomy (§4.1) is this problem restated for the LLM-coding-loop
setting, where the agent authors both code and partial oracles (tests).
Execution-based code-generation systems --- CodeT (Chen et al.~2022),
AlphaCodium (Ridnik et al.~2024) --- already combine generated tests
with execution filtering; yet by §4.1's collapse rule this is
same-source-by-cascade: the runtime executor is formal, but the
specification it executes (the test set) is agent-authored, so the
executor's blind spots collapse back to the model's. Self-consistency
(Wang et al.~2022) is same-source for the same reason --- diversity of
samples from one model does not diversify blind spots. Prior work thus
formalizes the agent's own proxy rather than decoupling the verification
source; our contribution is not the oracle concept but the explicit
two-axis organization (§3.3) and the salient/latent distinction that
separates what in-loop structural discipline can catch from what
requires an externalized oracle.

\textbf{LLM-as-judge and same-source verifier bias.} The LLM-as-judge
literature (Zheng et al.~2023) studies when and why an LLM evaluating
another LLM is reliable or biased --- including known
same-source/self-preference biases. Our self-certification paradox
(§3.1) is this bias axis applied to the coding-loop verification
problem: the same training distribution that produces the code produces
its validator. The contribution relative to LLM-as-judge is the
coding-loop-specific framing (proxy specification gaming via test
manipulation) and the Process/Outcome split as a schema-level
mitigation, not the discovery of same-source bias. The Evaluator Stress
Test (EST; Shihab et al.~2025) detects proxy gaming in RL/LLM-alignment
by separating exploitable-sensitivity from content-driven improvement
via controlled perturbations --- a score-decomposition axis distinct
from our integrity/adequacy (salient/latent) bifurcation (§4.3); its
external-perturbation audit is nonetheless a candidate realization of
the externalized heterogeneous oracle the latent regime demands.

\textbf{Industrial loop-engineering peers.} A 2026 cohort targets
adjacent ground. \emph{Spec Kit} (GitHub,
\url{https://github.com/github/spec-kit}) is spec-anchored with a
convergence primitive and a \texttt{constitution} anchor, but
single-ring and single-agent --- no adversarial reviewer, no
specification-gaming concept. \emph{Loki Mode}
(\url{https://github.com/asklokesh/loki-mode}) is the closest peer: a
tamper-evident proof artifact + drift detection + an Evidence Receipt
splitting Facts (deterministic) from Assessments (AI judgment) + a
multi-reviewer council --- a conceptual parallel to our Process/Outcome
split (§4.2) and heterogeneous adversarial ring (§4.4) --- same shape,
but weaker enforcement: Loki frames Facts/Assessments as a reporting
discipline rather than the schema-level invariant (an agent-barred enum)
that is the load-bearing property of §4.2, and we found no naming of
specification gaming as an orthogonal axis in its public docs.
\emph{zeroshot} (\url{https://github.com/the-open-engine/zeroshot}) adds
blind validation (validators cannot see the worker's context) with
fail-closed convergence across providers --- operationally close to our
heterogeneous subprocess. \emph{LoopX}
(\url{https://github.com/huangruiteng/loopx}) is a fourth peer of a
different shape: not a single-loop harness but a \textbf{cross-turn,
cross-runtime control plane} that externalizes objective, gates, todos,
evidence, and quota into an append-only event ledger while Codex, Claude
Code, or Cursor execute bounded turns. Its production verifiers are
same-runtime and, in its public docs, it does not name specification
gaming; on the orthogonal axis it is therefore silent, like the three
above. Two surfaces are nonetheless notable. First, its
benchmark/evaluation subsystem is a concrete \emph{deployed} instance of
the two-axis separation this paper prescribes (§5) --- an official,
externally-authored verifier scoring outcomes kept distinct from the
loop's own turn-readiness signal, with anti-gaming controls
(surrogate-run labelling, reward-leak prevention, and matched
treatment/baseline arms on the same case, source, model, reasoning
effort, timeout, and no-feedback budget). This instantiates the two-axis
\emph{evaluation separation} (the loop's turn-readiness metric kept
distinct from an externally-authored outcome metric), not a
specification-gaming injection-set methodology (ImpossibleBench, §8.3,
already fills that). Second, its \texttt{delivery\_outcome} enum --- a
structured self-report (\texttt{surface\_only} / \texttt{outcome\_gap} /
\texttt{outcome\_progress} / \texttt{primary\_goal\_outcome}) that
drives scheduling spend --- is a candidate \emph{salient} defense
against artifact-count Goodhart, but only partially so: the label is
executor-self-reported and validated only for a machine-visible
\emph{frontier delta} (form), not against the objective (substance),
leaving an irrelevant-change residue in the latent regime (§4.3).
Assessed against each project's public docs (absence claims bounded to
that surface), none of the four makes the irreducibility argument (§3.2)
or the two-axis frame (§3.3) explicit; the contribution of this paper is
to name the discipline they variously approximate without theorizing,
and to derive the platform-level structural guarantee (§4.4) that the
cohort approximates from different architectural positions ---
single-provider process-scoped harnesses from within, cross-runtime
control planes (LoopX) from a partial, seam-circumscribed position (see
the §4.4 note). LoopX additionally supplies the only observed deployed
instance of the two-axis separation among the cohort --- evidence that
the discipline is practised even where it is not theorized.

\textbf{Loop engineering as practice.} Osmani (2026) and concurrent
industry discussions (2026-06) named the \emph{practice} of designing
agent loops; they enumerate primitives (Automations, Worktrees, Skills,
Plugins, Sub-agents, Memory) and warn of unattended mistakes and
cognitive surrender. The practice framing is silent on the
orthogonal-axis question: it treats ``the loop runs and converges'' as
the goal, whereas this paper argues convergence (Axis A) is necessary
but insufficient, and that the verification source (Axis B) is the
independent discriminator.

\section{8. Open Problems}\label{open-problems}

\begin{enumerate}
\def\labelenumi{\arabic{enumi}.}
\tightlist
\item
  \textbf{Heterogeneous-oracle cost.} HITL annotation is not fully
  automatable; coverage is limited to annotated high-frequency cases ---
  novel failures regress to same-source verification.
\item
  \textbf{``Heterogeneous'' is itself PARTIAL.} Commercial model
  families share training data / RLHF overlap; cross-family review is
  only partial decoupling. Truly heterogeneous oracles (formal spec,
  production trace, human) remain the open frontier.
\item
  \textbf{From benchmark absence to defense-axis evaluation.} The
  empirical gap earlier drafts named --- no benchmark isolating
  specification gaming --- has closed on two fronts: observationally
  (SpecBench, Zhao et al.~2026, measures reward hacking via the
  visible/held-out pass-rate gap) and adversarially (ImpossibleBench,
  Zhong et al.~2025, injects direct conflicts between the
  natural-language specification and the unit tests on
  LiveCodeBench/SWE-bench, where any pass is a specification-violating
  shortcut; see also RHB, Thaman et al.~2026, on tool-using agents).
  What remains open is not benchmark construction but the
  \emph{defense-axis} evaluation this paper's two-axis frame predicts:
  does a heterogeneously decoupled oracle catch the specification gaming
  these benchmarks surface and a same-source judge misses?
\item
  \textbf{Formalization of ``verification-source decoupling degree.''}
  How to quantify Axis B as a measurable metric remains open --- a
  prerequisite for empirical comparison of systems on the second axis.
\item
  \textbf{Coverage of the root-cause path.} Acceptance criteria that
  cover only the symptom interface are bypassable by ``local
  satisfaction''; criteria that cover the upstream causal chain depend
  on diagnostic accuracy, which is itself a probabilistic hypothesis.
\item
  \textbf{Is process-scoped provider binding necessary or contingent?}
  (§4.4) A future harness with first-class multi-provider routing would
  need to re-derive the decoupling boundary by other means, or accept
  that the boundary becomes conventional.
\item
  \textbf{Heterogeneous-oracle reliability.} Positioning the
  heterogeneous oracle as the latent-regime defense leaves its own
  failure modes unexamined: false positives (correct code flagged as
  spec-gaming), oracle disagreement (two heterogeneous oracles returning
  conflicting verdicts), and error cascades (a wrong verdict triggering
  unnecessary rollback). The threshold at which oracle error outweighs
  defense value --- and how a false-positive rate interacts with the
  partial decoupling of §8.2 --- is open.
\item
  \textbf{Testability of the self-certification conjecture.} The
  necessity framing rests on the strong-form self-certification
  conjecture (§3.1), which we stipulate rather than establish.
  Empirically testing it --- e.g., measuring same-family-fresh-context
  vs cross-family verdict divergence on a known latent-proxy-gap test
  set, to determine whether same-family reviewers in fact share the
  verified's blind spots wholesale --- would either ground the
  ``logically necessary'' framing or weaken it to ``strongly advisable''
  (the weak-form fallback of §3.1).
\end{enumerate}

\section{References}\label{references}

\begin{itemize}
\tightlist
\item
  Krakovna, V. (2018). \emph{Specification gaming examples in AI.}
  \url{https://vkrakovna.wordpress.com/2018/04/02/specification-gaming-examples-in-ai/}
\item
  Krakovna, V. (2020). \emph{Specification gaming: the flip side of AI
  ingenuity.} DeepMind blog.
  \url{https://deepmind.google/blog/specification-gaming-the-flip-side-of-ai-ingenuity/}
\item
  Amodei, D., Olah, C., Steinhardt, J., Christiano, P., Schulman, J., \&
  Mané, D. (2016). \emph{Concrete Problems in AI Safety.}
  arXiv:1606.06565.
\item
  Pan, A., Bhatia, K., \& Steinhardt, J. (2022). \emph{The Effects of
  Reward Misspecification: Mapping and Mitigating Misaligned Models.}
  arXiv:2201.03544.
\item
  Manheim, D., \& Garrabrant, S. (2018). \emph{Categorizing Variants of
  Goodhart's Law.} arXiv:1803.04585.
\item
  Debenedetti, E., Shumailov, I., Fan, T., Hayes, J., Carlini, N.,
  Fabian, D., Kern, C., Shi, C., Terzis, A., \& Tramèr, F. (2025).
  \emph{Defeating Prompt Injections by Design (CaMeL).}
  arXiv:2503.18813.
\item
  Huang, D., Zhang, J. M., Luck, M., Bu, Q., Qing, Y., \& Cui, H.
  (2023). \emph{AgentCoder: Multi-Agent-based Code Generation with
  Iterative Testing and Optimisation.} arXiv:2312.13010.
\item
  Hong, S., Zhuge, M., Chen, J., et al.~(2023). \emph{MetaGPT: Meta
  Programming for a Multi-Agent Collaborative Framework.}
  arXiv:2308.00352.
\item
  Chan, C.-M., Chen, W., Su, Y., et al.~(2023). \emph{ChatEval: Towards
  Better LLM-based Evaluators through Multi-Agent Debate.}
  arXiv:2308.07201.
\item
  Liu, N. F., Lin, K., Hewitt, J., Paranjape, A., Bevilacqua, M.,
  Petroni, F., \& Liang, P. (2024). Lost in the Middle: How Language
  Models Use Long Contexts. \emph{Transactions of the Association for
  Computational Linguistics (TACL)}, 12, 157--173.
  \url{https://doi.org/10.1162/tacl_a_00638}
\item
  Osmani, A. (2026). \emph{Loop Engineering.}
  \url{https://addyosmani.com/blog/loop-engineering/}
\item
  Barr, E. T., Harman, M., McMinn, P., Shahbaz, M., \& Yoo, S. (2015).
  The Oracle Problem in Software Testing: A Survey. \emph{IEEE
  Transactions on Software Engineering}, 41(5), 507--525.
  \url{https://doi.org/10.1109/TSE.2014.2372785}
\item
  Chen, B., Zhang, F., Nguyen, A., Zan, D., Lin, Z., Lou, J.-G., \&
  Chen, W. (2022). \emph{CodeT: Code Generation with Generated Tests.}
  arXiv:2207.10397.
\item
  Ridnik, T., Kredo, D., \& Friedman, I. (2024). \emph{Code Generation
  with AlphaCodium: From Prompt Engineering to Flow Engineering.}
  arXiv:2401.08500.
\item
  Wang, X., Wei, J., Schuurmans, D., Le, Q., Chi, E., Narang, S.,
  Chowdhery, A., \& Zhou, D. (2022). Self-Consistency Improves Chain of
  Thought Reasoning in Language Models. \emph{ICLR 2023.}
  arXiv:2203.11171.
\item
  Zheng, L., Chiang, W.-L., Sheng, Y., et al.~(2023). \emph{Judging
  LLM-as-a-Judge with MT-Bench and Chatbot Arena.} arXiv:2306.05685.
\item
  Zhao, B., Srikanth, D., Wu, Y., \& Jiang, Z. (2026). \emph{SpecBench:
  Measuring Reward Hacking in Long-Horizon Coding Agents.}
  arXiv:2605.21384.
\item
  Shihab, I. F., Akter, S., \& Sharma, A. (2025). \emph{Detecting Proxy
  Gaming in RL and LLM Alignment via Evaluator Stress Tests.}
  arXiv:2507.05619. (ACL 2026 Findings.)
\item
  Zhong, Z., Raghunathan, A., \& Carlini, N. (2025).
  \emph{ImpossibleBench: Measuring LLMs' Propensity of Exploiting Test
  Cases.} arXiv:2510.20270.
\item
  Thaman, K. (2026). \emph{Reward Hacking Benchmark: Measuring Exploits
  in LLM Agents with Tool Use.} arXiv:2605.02964.
\item
  SolidForge --- reference implementation for §6. Repository:
  https://github.com/maskshell/solidforge
\end{itemize}

\end{document}
